\section*{Chapter \ref{ch:s2:systematic-macro} --- Systematic Macro}

\begin{solution}{exo:s2:systematic-macro:1}
$(0.84 + 0.78 + 0.33 + 0.29)/\sqrt4 = 1.12$.
\end{solution}

\begin{solution}{exo:s2:systematic-macro:2}
$(0.84 + 0.78)/\sqrt{2 + 2 \times (-0.43)} = 1.52$: a negative correlation raises the combination's Sharpe ratio above the sum over $\sqrt2$.
\end{solution}

\begin{solution}{exo:s2:systematic-macro:3}
$1 - 0.5^{1/3} = 21\%$.
\end{solution}

\begin{solution}{exo:s2:systematic-macro:4}
The trend lasts about a year, and the five-year price move still contains the recent part of it: minus that move is short the markets that trended up,
the opposite of momentum. The anchor-based signal carries no such overlap.
\end{solution}

\begin{solution}{exo:s2:systematic-macro:5}
Macro data are revised, sometimes heavily; a backtest on revised data uses numbers no one had at the time and overstates the signal.
\end{solution}

\begin{solution}{exo:s2:systematic-macro:6}
Asset classes have different volatilities and common factors; within a class the book is neutral to the class's own direction and trades only the
signal, and each class's risk can be budgeted separately.
\end{solution}

\begin{solution}{exo:s2:systematic-macro:7}
Value falls from 0.84 to 0.69 and the combination from 1.15 to 1.05: a noisier anchor dilutes value, and the other three signals carry the book.
\end{solution}

\begin{solution}{exo:s2:systematic-macro:8}
The Sharpe ratio comes from a planted market and from signals whose correlations may rise together in a crisis; at 30\% volatility the combination's
worst drawdown of 21\% at 10\% volatility would be about 63\%. Size on stressed correlations and drawdowns, not on the measured ratio.
\end{solution}

\begin{solution}{pb:s2:systematic-macro:1}
\begin{enumerate}
\item Price against a fundamental anchor; purchasing-power parity, real yields, earnings yields.
\item Long improving macro trends, short deteriorating; data releases against forecasts.
\item Long-short by rank within each class, monthly, scaled to 10\% volatility.
\item Value and momentum premia everywhere, negatively correlated, with common factors across asset classes.
\item A gap of 10\% sd reverting over three years; a surprise index reverting over 60 days.
\item The gap through an anchor with a persistent 10\% error; the surprise with noise of half its size.
\item It dilutes the signal: at 30\% error value earns 0.69.
\item \omterm{def:s2:systematic-macro:momentum}{Macro momentum}: long improving fundamentals, short deteriorating, with low correlation to asset classes.
\item Value 0.84, reversal $-0.84$, momentum 0.78, carry 0.33, surprise 0.29.
\item Near zero except reversal and momentum, $-0.43$.
\item 1.15 with anchor value, 0.31 with price value, 0.81 without value.
\item It overlaps with the trend and becomes its mirror.
\item \textbf{Named result.} Value 0.84, momentum 0.78, carry 0.33 and surprise 0.29, nearly uncorrelated, combine to 1.15 with a smaller drawdown
than any; value from prices is momentum's mirror ($-0.43$ correlation) and drags the combination to 0.31.
\item The sum of Sharpe ratios over the square root of their number: 1.12.
\item Correlations tend to rise; the combination's benefit shrinks when needed most.
\item Equal to start, then from trailing correlations and volatilities, capped.
\item Point-in-time data, signals fixed in advance, costs, several seeds or periods.
\item The economic-surprise tilt.
\item It adds value and macro data to Book~8's trend and carry.
\item Each is paid for a different risk, and the risks do not arrive together.
\end{enumerate}
\end{solution}

\begin{solution}{iq:s2:systematic-macro:1}
Currencies: real exchange rates against purchasing-power parity; bonds: real yields or term premia against history; equity indices: earnings or book
yields against bond yields or their own history.
\end{solution}

\begin{solution}{iq:s2:systematic-macro:2}
Momentum buys what has risen, value what has fallen relative to its anchor, so they often take opposite sides; the negative correlation makes the
combination steadier than either.
\end{solution}

\begin{solution}{iq:s2:systematic-macro:3}
Net them as the book's risk budgets say; the signals are designed to disagree at times, and overriding one because the other exists breaks the
combination's diversification.
\end{solution}

\begin{solution}{iq:s2:systematic-macro:4}
A trend reversal across all markets, a carry crash, a correlation spike across asset classes, and a data shock (a large surprise) that moves all
countries together.
\end{solution}

\begin{solution}{iq:s2:systematic-macro:5}
Store every release with its date and time, its first print and each revision, and the consensus forecast before it; query by the knowledge date, not
the reference period.
\end{solution}

\begin{solution}{iq:s2:systematic-macro:6}
Each signal $i$ is scaled to volatility $\sigma$ and earns $S_i\sigma$; the equal-weight sum earns $\sigma\sum S_i$ with variance $\sigma^2(n +
n(n - 1)\rho)$, so the ratio is $\sum_i S_i/\sqrt{n + n(n - 1)\rho}$.
\end{solution}
